% Generated by: tools/render_public_repo_scores.py
% Source: paper/data/public-repo-scores-20260505T184426Z.json
% Command: python3 tools/render_public_repo_scores.py --source paper/data/public-repo-scores-20260505T184426Z.json --out paper/tex/generated/public_repo_score_tables.tex
% DO NOT EDIT BY HAND.
% Run root: /home/ubuntu/jankscore/runs/20260505T184426Z
% Generated at: 2026-05-05T18:45:33Z
% Jankurai version: 0.7.0

\newcommand{\PublicRepoScoreTopTable}{%
\begin{table*}[t]
\caption{Top observed advisory public-repository scores. The top score is best observed in this run, not a certification threshold.}
\label{tab:public-repo-top-scores}
\centering
\scriptsize
\setlength{\tabcolsep}{3pt}
\begin{tabularx}{\textwidth}{R{0.045\textwidth} L{0.245\textwidth} R{0.105\textwidth} R{0.075\textwidth} R{0.065\textwidth} Y}
\toprule
\JKTableHeader
\textbf{Rank} & \textbf{Repo} & \textbf{Score} & \textbf{Findings} & \textbf{Hard} & \textbf{Dominant gaps} \\
\midrule
\JKDenseRows
1 & \path{zed-industries/zed} & \JKScoreCellBest{48} & 2,029 & 2,016 & code shape/semantic surface; Python/\allowbreak{}polyglot hygiene; audit CI replacement \\
2 & \path{astral-sh/ruff} & \JKScoreCell{45} & 2,752 & 2,739 & code shape/semantic surface; Python/\allowbreak{}polyglot hygiene; audit CI replacement \\
3 & \path{denoland/deno} & \JKScoreCell{43} & 1,615 & 1,601 & code shape/semantic surface; audit CI replacement; context routing \\
4 & \path{RightNow-AI/openfang} & \JKScoreCell{42} & 690 & 677 & code shape/semantic surface; Python/\allowbreak{}polyglot hygiene; audit CI replacement \\
5 & \path{meilisearch/meilisearch} & \JKScoreCell{42} & 696 & 685 & code shape/semantic surface; audit CI replacement; proof routing \\
6 & \path{nearai/ironclaw} & \JKScoreCell{42} & 1,798 & 1,785 & code shape/semantic surface; Python/\allowbreak{}polyglot hygiene; audit CI replacement \\
7 & \path{googleworkspace/cli} & \JKScoreCell{41} & 97 & 86 & code shape/semantic surface; audit CI replacement; proof routing \\
8 & \path{zeroclaw-labs/zeroclaw} & \JKScoreCell{40} & 1,213 & 1,200 & code shape/semantic surface; Python/\allowbreak{}polyglot hygiene; audit CI replacement \\
9 & \path{kvnxiao/tauri-tanstack-start-react-template} & \JKScoreCell{39} & 22 & 10 & audit CI replacement; context routing; proof routing \\
10 & \path{vercel-labs/agent-browser} & \JKScoreCell{38} & 189 & 178 & code shape/semantic surface; audit CI replacement; security/supply chain \\
\bottomrule
\end{tabularx}
\JKResetRows
\end{table*}
}

\newcommand{\PublicRepoScoreAggregateTable}{%
\begin{table}[t]
\caption{Aggregate advisory scan posture.}
\label{tab:public-repo-aggregate}
\centering
\scriptsize
\setlength{\tabcolsep}{4pt}
\begin{tabularx}{\columnwidth}{L{0.60\columnwidth} R{0.28\columnwidth}}
\toprule
\JKTableHeader
\textbf{Metric} & \textbf{Value} \\
\midrule
\JKDenseRows
Repositories scanned & 30 \\
Successful scans & 30 \\
Failed scans & 0 \\
Minimum score & 14 \\
Maximum score & 48 (best observed) \\
Average score & 33.4 \\
Median score (upper middle) & 34 \\
Findings total & 15,391 \\
Hard findings total & 15,017 \\
Hard findings share & 97.6\% \\
\bottomrule
\end{tabularx}
\JKResetRows
\end{table}
}

\newcommand{\PublicRepoWeakDimensionTable}{%
\begin{table}[t]
\caption{Most frequent weak-dimension sets in the advisory scan.}
\label{tab:public-repo-weak-dimensions}
\centering
\scriptsize
\setlength{\tabcolsep}{4pt}
\begin{tabularx}{\columnwidth}{Y R{0.18\columnwidth}}
\toprule
\JKTableHeader
\textbf{Weak dimension} & \textbf{Repos} \\
\midrule
\JKDenseRows
Jankurai tool adoption and CI replacement & 29/30 \\
Code shape and semantic surface & 28/30 \\
Context economy and agent instructions & 12/30 \\
Python containment and polyglot hygiene & 11/30 \\
Proof lanes and test routing & 8/30 \\
\bottomrule
\end{tabularx}
\JKResetRows
\end{table}
}

\newcommand{\PublicRepoScoreAppendixTable}{%
\begingroup
\scriptsize
\setlength{\tabcolsep}{3pt}
\setlength{\LTleft}{0pt}
\setlength{\LTright}{0pt}
\JKDenseRows
\begin{longtable}{@{}R{0.040\textwidth} R{0.045\textwidth} L{0.235\textwidth} R{0.055\textwidth} R{0.070\textwidth} R{0.065\textwidth} L{0.355\textwidth}@{}}
\caption{Full 30-repository advisory scoring run.}\label{tab:public-repo-full}\\
\toprule
\JKTableHeader
\textbf{Rank} & \textbf{Run \#} & \textbf{Repo} & \textbf{Score} & \textbf{Findings} & \textbf{Hard} & \textbf{Dominant gaps} \\
\midrule
\endfirsthead
\toprule
\JKTableHeader
\textbf{Rank} & \textbf{Run \#} & \textbf{Repo} & \textbf{Score} & \textbf{Findings} & \textbf{Hard} & \textbf{Dominant gaps} \\
\midrule
\endhead
1 & 25 & \path{zed-industries/zed} & \JKScoreCellBest{48} & 2,029 & 2,016 & code shape/semantic surface; Python/\allowbreak{}polyglot hygiene; audit CI replacement \\
2 & 26 & \path{astral-sh/ruff} & \JKScoreCell{45} & 2,752 & 2,739 & code shape/semantic surface; Python/\allowbreak{}polyglot hygiene; audit CI replacement \\
3 & 23 & \path{denoland/deno} & \JKScoreCell{43} & 1,615 & 1,601 & code shape/semantic surface; audit CI replacement; context routing \\
4 & 6 & \path{RightNow-AI/openfang} & \JKScoreCell{42} & 690 & 677 & code shape/semantic surface; Python/\allowbreak{}polyglot hygiene; audit CI replacement \\
5 & 29 & \path{meilisearch/meilisearch} & \JKScoreCell{42} & 696 & 685 & code shape/semantic surface; audit CI replacement; proof routing \\
6 & 9 & \path{nearai/ironclaw} & \JKScoreCell{42} & 1,798 & 1,785 & code shape/semantic surface; Python/\allowbreak{}polyglot hygiene; audit CI replacement \\
7 & 4 & \path{googleworkspace/cli} & \JKScoreCell{41} & 97 & 86 & code shape/semantic surface; audit CI replacement; proof routing \\
8 & 3 & \path{zeroclaw-labs/zeroclaw} & \JKScoreCell{40} & 1,213 & 1,200 & code shape/semantic surface; Python/\allowbreak{}polyglot hygiene; audit CI replacement \\
9 & 12 & \path{kvnxiao/tauri-tanstack-start-react-template} & \JKScoreCell{39} & 22 & 10 & audit CI replacement; context routing; proof routing \\
10 & 2 & \path{vercel-labs/agent-browser} & \JKScoreCell{38} & 189 & 178 & code shape/semantic surface; audit CI replacement; security/supply chain \\
11 & 21 & \path{tauri-apps/tauri} & \JKScoreCell{38} & 304 & 292 & code shape/semantic surface; audit CI replacement; context routing \\
12 & 24 & \path{astral-sh/uv} & \JKScoreCell{38} & 1,140 & 1,128 & code shape/semantic surface; Python/\allowbreak{}polyglot hygiene; audit CI replacement \\
13 & 19 & \path{kevinpiao1025/tauri-react-typescript-tailwind} & \JKScoreCell{36} & 22 & 10 & audit CI replacement; proof routing; context routing \\
14 & 5 & \path{AlexsJones/llmfit} & \JKScoreCell{35} & 107 & 94 & code shape/semantic surface; Python/\allowbreak{}polyglot hygiene; audit CI replacement \\
15 & 27 & \path{alacritty/alacritty} & \JKScoreCell{34} & 110 & 98 & code shape/semantic surface; context routing; audit CI replacement \\
16 & 16 & \path{exit-zero-labs/threat-forge} & \JKScoreCell{32} & 134 & 121 & code shape/semantic surface; audit CI replacement; context routing \\
17 & 8 & \path{microsoft/RustTraining} & \JKScoreCell{31} & 34 & 21 & code shape/semantic surface; audit CI replacement; context routing \\
18 & 28 & \path{BurntSushi/ripgrep} & \JKScoreCell{31} & 72 & 60 & code shape/semantic surface; audit CI replacement; context routing \\
19 & 11 & \path{octasoft-ltd/wsl-ui} & \JKScoreCell{31} & 314 & 302 & code shape/semantic surface; audit CI replacement; context routing \\
20 & 30 & \path{typst/typst} & \JKScoreCell{31} & 382 & 370 & code shape/semantic surface; audit CI replacement; context routing \\
21 & 22 & \path{rustdesk/rustdesk} & \JKScoreCell{30} & 648 & 636 & code shape/semantic surface; proof routing; Python/\allowbreak{}polyglot hygiene \\
22 & 1 & \path{rtk-ai/rtk} & \JKScoreCell{28} & 397 & 383 & code shape/semantic surface; audit CI replacement; context routing \\
23 & 20 & \path{MarkShawn2020/lovtauri} & \JKScoreCell{26} & 30 & 18 & audit CI replacement; proof routing; code shape/semantic surface \\
24 & 18 & \path{Duri686/RustQuantLab} & \JKScoreCell{26} & 42 & 30 & code shape/semantic surface; audit CI replacement; context routing \\
25 & 10 & \path{h4ckf0r0day/obscura} & \JKScoreCell{26} & 60 & 48 & code shape/semantic surface; audit CI replacement; proof routing \\
26 & 7 & \path{xai-org/x-algorithm} & \JKScoreCell{24} & 38 & 25 & code shape/semantic surface; audit CI replacement; Python/\allowbreak{}polyglot hygiene \\
27 & 14 & \path{sergioadevita/notemac-plus-plus} & \JKScoreCell{24} & 293 & 279 & code shape/semantic surface; audit CI replacement; Python/\allowbreak{}polyglot hygiene \\
28 & 17 & \path{lostf1sh/rustune} & \JKScoreCell{23} & 36 & 24 & code shape/semantic surface; audit CI replacement; proof routing \\
29 & 13 & \path{fudanglp/tauri-fastapi-full-stack-template} & \JKScoreCell{23} & 69 & 56 & code shape/semantic surface; audit CI replacement; Python/\allowbreak{}polyglot hygiene \\
30 & 15 & \path{ianho7/maptoposter-online} & \JKScoreCell{14} & 58 & 45 & build speed; code shape/semantic surface; audit CI replacement \\
\bottomrule
\end{longtable}
\JKResetRows
\endgroup
}


\section{Merge-Witness Readiness Scan}

The strongest field evidence in this edition is deliberately advisory. On May 5, 2026, Jankurai 0.7.0 scanned 30 public GitHub repositories against the current target profile. The run succeeded on all 30 repositories and produced local JSON, Markdown, clone, and log artifacts under one run root \cite{jankuraiPublicRepoRun2026}. This is not a certification result, not a defect claim, and not an empirical incident study. It is a posture scan asking whether prominent repositories expose merge-witness-ready ownership, proof routing, generated-zone receipts, context routing, proof freshness, and repair evidence.

The sample includes mature and widely used public projects such as Zed, Ruff, Deno, Meilisearch, Tauri, and ripgrep \cite{zedGithub2026,ruffGithub2026,denoGithub2026,meilisearchGithub2026,tauriGithub2026,ripgrepGithub2026}. Many of these projects have serious engineering cultures, active maintainers, and conventional CI. The surprising result is therefore not that small or immature repositories score low. It is that conventional maturity and popularity do not automatically produce the repository-alignment surfaces Jankurai expects.

The top observed score was 48, and every scanned repository remained below the 50-point red band. The average score was 33.4. Across the run, Jankurai emitted 15,391 findings, including 15,017 hard findings. Those counts should be read as evidence gaps and repair opportunities under Jankurai policy, not as accusations against maintainers.

\PublicRepoScoreAggregateTable

The common missing surfaces were consistent with the paper's thesis. Repositories often had tests and CI, but did not expose an audit lane as a CI replacement artifact; did not bind owner maps, test maps, and generated-zone policy into a merge witness; did not leave fresh proof receipts for changed paths; and did not serialize repair queues in a form that a weaker implementation agent could consume. This matches the broader AI-assisted development evidence: DORA frames AI as an amplifier of existing delivery capability, not an automatic substitute for engineering discipline; security and secret-sprawl reports show plausible generated work can increase evidence risk; developer surveys and agent benchmarks show that nearly-correct output, environment setup, and context retrieval remain expensive failure modes \cite{doraStateAIAssistedSoftwareDevelopment2025,veracodeGenaiCodeSecurity2025,gitguardianSecretsSprawl2026,gitguardianSecretsSprawl2026Blog,stackOverflowAISurvey2025,sweAgentAci2024,setupBench2025,contextBench2026}.

\PublicRepoWeakDimensionTable

The weak-dimension frequency makes the gap concrete. Jankurai tool adoption and CI replacement appeared in 29 of 30 weak-dimension sets. Code shape and semantic surface appeared in 28 of 30. Context economy and agent instructions appeared in 12 of 30, Python/polyglot hygiene in 11 of 30, and proof lanes and test routing in 8 of 30. These are not claims that the projects lack quality. They are claims that the repositories generally do not publish the specific merge-evidence interface required for a Jankurai witness.

The caveat matters. HLT marker counts, hard caps, and target-profile assumptions are strict policy proxies. A task marker, large file, Python placement, generated-boundary heuristic, or missing Jankurai CI artifact can be a false positive or a locally accepted tradeoff. The advisory scan is useful because it names repairable evidence gaps, not because it proves insecurity or nonconformance. Popularity, language choice, and conventional CI do not by themselves create repository alignment. The missing merge witness is exactly the standard gap Jankurai targets.

The full named table remains in Appendix~F. Keeping the main body aggregate-first avoids turning an exploratory readiness scan into a public quality leaderboard.
